% ---
% id: meat-core/prop:optol-nobound
% kind: prop
% title: "Instability of the rate under measurement perturbations"
% summary: >-
%   Two exactly admissible measurements on Q within ϵ of each other can have
%   rate bounds α/(α−1)·N bits per round apart, for every positive integer
%   N ≤ 53 + log₂ϵ (up to 52). So closeness of the declared measurement
%   cannot be relied on to carry a key-rate claim to a nearby declaration.
% prominence: standard
% status: accepted
% uses: [def:feasible, def:declaration, prop:value-range, thm:keylength, def:max-key-length, prop:trivial-alice-system, foundations/def:kl, def:admissible-declaration]
% mentions: [rem:representability-ceiling]
% assumes: []
% refs: []
% provenance: newly-derived
% lean:
% notes: "The statement is over declarable binary64 families, not real-valued ones, so its gap, between families of finite value (g_0 > 0), is capped by the number model (N at most 52); the exactly admissible family (0, 1) has h_low = +infinity and carries no key-length claim (req:infinite-value), so it is set aside; do not restate the finite gap as unbounded. Epsilon is any number in (0, 1/2] with 2^N <= epsilon 2^53, and Step 2 builds the families at the largest power of two epsilon' <= epsilon: the construction needs epsilon' 2^-N to be a multiple of 2^-53, which a general dyadic epsilon does not give (epsilon = 2^-2 (1 + 2^-52), N = 1). Part (i) is cited by number from rem:classes."
% ---

One might admit a declaration whose measurement operators are only close
to exact ones, with an error allowance depending on the distance. The
following proposition shows that no allowance computed from the declared
dimensions, alphabet sizes and $\alpha$ can be useful. It compares values
$\hlow$ of \eqref{eq:hlow-def}, the infimum of the objective
\eqref{eq:objective} over the acceptance set and the feasible states. That
objective adds to its entropy terms the acceptance-test penalty
$\frac{\alpha}{\alpha-1}\sfD(\mathbf q\|\Phimap[\omega])$, where
$\mathbf q$ is a frequency distribution from the acceptance set and
$\Phimap[\omega]$ the distribution of the test statistic under $\omega$;
the penalty grows without bound as $\mathbf q$ demands an outcome that
$\omega$ makes improbable. The hypothesis $2^{N}\le\epsilon\,2^{53}$ is
what the binary64 number model allows
(Remark~\ref{rem:representability-ceiling}). Part~(iii) is arithmetic on
the formula $\ell_{\max}$ of Definition~\ref{def:max-key-length}, which
Theorem~\ref{thm:keylength} reads as a key length on $(1,2)$ under its own
assumptions.

\begin{proposition}\label{prop:optol-nobound}
Let $\alpha\in(1,\infty)$ and $B>0$, write $N\coloneqq\lceil B\rceil$, and
let $\epsilon\in(0,\tfrac12]$ satisfy $2^{N}\le\epsilon\,2^{53}$.
Write $\Decl^{(i)}$ for the declaration obtained from a given one by
replacing its measurement $\{F_v\}$ on $\Qsys$ by $\{F^{(i)}_v\}$, and
$\hlow[\,\cdot\,]$ for the value \eqref{eq:hlow-def} of a declaration.
\begin{enumerate}[label=(\roman*),leftmargin=2.4em]
\item There is a declaration $\Decl$ satisfying Definition~\ref{def:admissible-declaration}
in full, with $\{F_v\}$ itself \tref{exactly admissible}, and there are two
exactly admissible families $\{F^{(1)}_v\}$, $\{F^{(2)}_v\}$ with
$\max_{v}\|F^{(i)}_v-F_v\|\le\epsilon$ for $i=1,2$, such that
\begin{equation}\label{eq:optol-gap}
  \hlow\big[\Decl^{(1)}\big]-\hlow\big[\Decl^{(2)}\big]\;\ge\;B .
\end{equation}
The dimensions and alphabets of $\Decl$ are $d_A=d_{\Qsys}=1$, one
sector, $|\Xset|=1$, $|\Vset|=|\Chatset|=2$ and $d_Z=1$, the same for
every $B$ and every $\epsilon$.
\item Let $\mathrm{err}$ be a function of the declared dimensions,
alphabet sizes and $\alpha$ whose value on the data of (i) is less than
$\frac{\alpha}{\alpha-1}\,N$. Then the bound
$\big|\hlow[\Decl^{(1)}]-\hlow[\Decl^{(2)}]\big|\le\mathrm{err}$, with
$\mathrm{err}$ evaluated on the data of the declaration, fails for some
declaration and some pair of exactly admissible families within
$\epsilon$ of each other.
\item For $\alpha\in(1,2)$ and every number $n$ of rounds, the
\tref[maximal key length]{maximal admissible key lengths} $\ell_{\max}$ of
$\Decl^{(1)}$ and $\Decl^{(2)}$
(Definition~\ref{def:max-key-length}), computed with the same $\lEC$, $\eev$ and
$\epa$, differ by at least $nB-1$ bits whenever the smaller of the two is
positive.
\end{enumerate}
\end{proposition}

\begin{proof}
\emph{Idea.} The instance has one-dimensional systems, so a single state;
the measurement on $\Qsys$ has two outcomes, $0$ with probability $g_0$
and $1$ with probability $g_1$; the acceptance test demands outcome~$0$ in
every round; and no round contributes key, so no acceptance value is
key-live. Then $\hlow$ is the acceptance-test penalty alone,
$\frac{\alpha}{\alpha-1}\log_2\frac1{g_0}$. Dividing $g_0$, the largest
power of two at most $\epsilon$, by $2^{N}$ raises $\hlow$ by
$\frac{\alpha}{\alpha-1}N$ and moves each entry of the family by at most
$\epsilon$, so the declared family and its perturbation lie within
$\epsilon$ of each other.

\emph{Step 1 (the instance).} Take the instance of
Proposition~\ref{prop:value-range}\,(iii): $d_A=d_{\Qsys}=1$ (so both source modes
coincide, Proposition~\ref{prop:trivial-alice-system}, and $\Feas=\{[1]\}$), one
sector, $\Xset=\{0\}$ with $M_0=[1]$, $\Vset=\{0,1\}$, $a_A$ constant and
$a_B=\idmap$, $\phi(\kappa)=c_B$ so that $\Chatset=\Vset$, all-fail
sifting ($\varsigma\equiv\mathrm{fail}$, $r\equiv\perp$, $\hold=\mathsf B$
with $r_{\mathsf B}\equiv\perp$, hence $\Chkey=\varnothing$),
$\CList=\varnothing$, the default \tref{party partition}, and the singleton
acceptance set $\acSet=\{(1,0)\}$. For scalars $g_0,g_1\ge0$ write
$\Decl_{(g_0,g_1)}$ for the declaration whose measurement on $\Qsys$ is
$F_0=[g_0]$, $F_1=[g_1]$. Every clause of Definition~\ref{def:admissible-declaration} other
than the $\{F_v\}$-clause of (iii) holds for it: all systems and
alphabets are finite, sifting is announcement-determined, the key map is
party-computable with $r_{\mathsf B}\equiv\perp$,
$\Feas=\{[1]\}\ne\varnothing$, $\acSet$ is closed and convex, and
$\dim\Hs{\Etil}\ge1$. That clause holds exactly when $g_0+g_1=1$, and the
family is exactly admissible in the sense of
Definition~\ref{def:declaration} exactly when in addition both entries are
\tref{binary64 numbers} and the one addition of \eqref{eq:declorder} is exact.
Since $\Chkey=\varnothing$, the objective \eqref{eq:objective} keeps only
its first sum, the acceptance-test penalty
$\frac{\alpha}{\alpha-1}\sum_{\hat c}q(\hat c)\log_2\frac{q(\hat c)}{P_{\hat c}(\omega)}$,
and $\hlow$ of \eqref{eq:hlow-def} is its infimum over
$\mathbf q\in\acSet=\{(1,0)\}$ and $\omega\in\Feas$. $\Feas$ is the single
state $[1]$, at which $P_{\hat c_0}=g_0$ and $P_{\hat c_1}=g_1$; with the
convention $0\log_2\frac0v=0$ of \xref{foundations/def:kl} this gives
\begin{equation}\label{eq:optol-instance}
  \hlow\big[\Decl_{(g_0,g_1)}\big]
  =\frac{\alpha}{\alpha-1}\Big(1\cdot\log_2\frac1{g_0}
  +0\cdot\log_2\frac{0}{g_1}\Big)
  =\frac{\alpha}{\alpha-1}\,\log_2\frac1{g_0}
  \qquad (g_0>0,\ g_1\ge0),
\end{equation}
a computation that uses no property of $(g_0,g_1)$ beyond those two
inequalities, in particular not $g_0+g_1=1$.

\emph{Step 2 (part (i)).} Let $\epsilon'$ be the largest power of two
with $\epsilon'\le\epsilon$, so $\epsilon'=2^{-k}$ for an integer
$k\ge1$. The hypothesis says $2^{N-53}\le\epsilon$, and $2^{N-53}$ is a
power of two, so $2^{N-53}\le\epsilon'$ as well. Take
$\Decl\coloneqq\Decl_{(\epsilon',\,1-\epsilon')}$,
$\{F^{(2)}_v\}\coloneqq\{F_v\}$ (distance $0$) and $\{F^{(1)}_v\}$ the
POVM of $\Decl_{(\epsilon'2^{-N},\,1-\epsilon'2^{-N})}$, whose distance to
$\{F_v\}$ is $\epsilon'(1-2^{-N})\le\epsilon'\le\epsilon$ in each of its
two entries. All three families have non-negative entries summing to $1$;
each has $|\Vset|=2$, so \eqref{eq:declorder} performs one addition. By
$2^{N-53}\le\epsilon'$, $\epsilon'2^{-N}=2^{-k-N}$ is at least $2^{-53}$,
so $\epsilon'2^{-N}$ and $\epsilon'$ are powers of two in
$[2^{-53},\tfrac12]$, hence binary64 numbers and multiples of $2^{-53}$;
$\epsilon=\tfrac12$ with $N\le52$ is one such choice. Then every entry is
exactly representable ($1-\epsilon'$ and $1-\epsilon'2^{-N}$ are
multiples of $2^{-53}$ in $[\tfrac12,1)$, where consecutive binary64
numbers are $2^{-53}$ apart), and the single addition returns the
representable $1$ and is therefore exact, which is the clause of
Definition~\ref{def:declaration} that no addition of \eqref{eq:declorder}
rounds. So all three families are exactly admissible and $\Decl$
satisfies Definition~\ref{def:admissible-declaration} in full. By \eqref{eq:optol-instance},
\begin{equation*}
  \hlow\big[\Decl^{(1)}\big]-\hlow\big[\Decl^{(2)}\big]
  =\frac{\alpha}{\alpha-1}
  \Big(\log_2\frac{1}{\epsilon'2^{-N}}-\log_2\frac1{\epsilon'}\Big)
  =\frac{\alpha}{\alpha-1}\,N\;\ge\;N\;\ge\;B ,
\end{equation*}
using $\frac{\alpha}{\alpha-1}>1$ for $\alpha>1$.

\emph{Step 3 (part (ii)).} Let $\mathrm{err}$ be as in (ii). The
declaration $\Decl$ of Step~2 has the dimensions and alphabets listed in
(i), which are the same for every $B$ and $\epsilon$, so $\mathrm{err}$
takes one value on its data, and that value is less than
$\frac{\alpha}{\alpha-1}\,N$. The families $\{F^{(1)}_v\}$ and
$\{F^{(2)}_v\}=\{F_v\}$ of Step~2 are exactly admissible and lie within
$\epsilon$ of each other, and their values differ by
$\frac{\alpha}{\alpha-1}\,N$, which exceeds the prescribed $B$ for every
$B$ the representability condition admits and exceeds that value of
$\mathrm{err}$. So the bound fails for this pair. Since the value of
$\mathrm{err}$ does not depend on $B$, the argument may be run at the
largest $N$ the representability condition admits.

\emph{Step 4 (part (iii)).} Let $\alpha\in(1,2)$. By
Definition~\ref{def:max-key-length}, for $n$ rounds
\begin{equation*}
  \ell_{\max}=\max\Big\{0,\Big\lfloor n\hlow-\lEC
  -\Big\lceil\log_2\frac1{\eev}\Big\rceil
  -\frac{\alpha}{\alpha-1}\log_2\frac1{\epa}+1\Big\rfloor\Big\};
\end{equation*}
the terms after $n\hlow$ do not depend on $\hlow$ and are the same for
both declarations. Insert the gap \eqref{eq:optol-gap}: if the smaller of the
two lengths is positive, the two lengths differ by at least $nB-1$ bits.
\end{proof}
